% Page 020

\seite{20}

\jahresueberschrift{Das Jahr 1892}
\pstart
Die Schülerzahl betrug in diesem Jahre 50.\ob
Im Winter dieses Jahres gerieth der Gemeinderath mit dem\ob
größten Theile der Bürger in Konflikt, weil derselbe den\ob
Beschluß faßte, ein Drittel des Holzlooses den Berechtigten\ob
zu entziehen. Von den 46 Berechtigten des Dorfes bezog\ob
\margmark{Der Waldreichthum\\Sefferweichs.}%
jeder gewöhnlich jährlich 6 Raummeter Buchenscheitholz\ob
nebst Reiser und zwei Raummeter Tannenholz. 32\ob
Bürger legten Berufung ein bei dem zuständigen\ob
Kreisausschusse, jedoch ohne Erfolg. Das Holz kam zur\ob
Versteigerung.

Die Gemeinde Sefferweich hat ungefähr 800 Morgen Wald,\ob
davon sind 300 Morgen Nieder- und 500 M.\ Hochwald.\ob
Letzterer liegt auf dem eingangs erwähnten\ob
Siebengemeinden Walde. In früherer Zeit konnte die\ob
Gemeinde sämmtliche Ländereien im Distrikt Klaftbrunnen\ob
erhalten, 1500 Morgen, wenn sie auf ihren Antheil\ob
am Siebengemeinden Walde verzichtete. Die Gemeinde\ob
\margmark{Die frühere Ka-\\pelle auf dem\\Klaftbrunnen.}%
ging nicht darauf ein, wiewohl dies nach Ansicht\ob
von Sachkundigen ein vortheilhafter Tausch\ob
gewesen wäre.

Auf dem Klaftbrunnen befand sich früher eine kleine\ob
Kapelle, die jetzt in einen Stall umgewandelt ist. Über\ob
der Kapelle war die Wohnung des Herrn Pfarrer Dunckel\unsicher, der\ob
hier alltäglich die hl. Messe las. Die Bauart der Kapelle ist\ob
heute noch ersichtlich. (Heutiger Besitzer: Raths Johannes\unsicher)
\pend
